\section{Summary of Findings}

The main aim of this thesis was to determine the effectiveness of cattle visual representations, in addition to multi-task learning, for monitoring dairy cattle. The tasks carried out under the thesis were: body condition scoring(BCS), behaviour recognition and re-identification of the cattle. The entire procedure took place through the process of record keeping, leakage-aware evaluation protocols, and matched-population controls to make sure its reproducible and reliable.

In the single-task experiments, it was possible to identify obvious advantages of cattle-oriented representations, but the degree of improvement was not equal across tasks and metrics.:
\begin{itemize}
    \item As for BCS, there were 7,549 images in ScienceDB test set as detection and segmentation operated well on them, which constituted 93.89\% of the whole test set. Cattle image cropping and binary masks yielded better performance in comparison with RGB representation. The value of MAE was reduced from 0.1929 to 0.1709 BCS units; Acc@1 within $\pm$0.25 was increased from 84.95\% to 89.40\%, however, exact classification was 43.57\%. At the same time, balanced accuracy and Macro-F1 decreased slightly, from 40.23\% to 39.70\% and from 0.4074 to 0.4039, correspondingly.
    \item As for Behavior task, tests were conducted using 780 video clips consisting of 8 frames per clip that constitute 96.42\% of matching data. Cropped cattle images, mask in the foreground, and small 1D temporal network were found to be beneficial in several ways compared to the single frame RGB model. Balanced accuracy increased to 74.43\% from 71.30\%, test loss decreased to 0.4430 from 0.5505, and walking F1 increased to 0.2456 from 0.2174. Nevertheless, the overall accuracy reduced to 87.44\% from 88.46\%. The approach was more effective for the feedlot camera images but somewhat less effective for the CVB surveillance images.
    \item In the case of Cattle Re-Identification, the model was evaluated on 69 unobserved cattle using the SideView Protocol A. Utilization of the ground-truth masks for creating target-centered crops enhanced first-match accuracy for both queries. In the case of Barn-to-Parlor, Rank-1 increased from 58.64\% to 63.90\%, while mAP increased from 38.32\% to 40.68\%. In the case of Snapshot-to-Parlor, Rank-1 increased from 38.88\% to 62.93\%, while mAP increased from 27.05\% to 40.42\%. The substantial gain in the snapshot performance indicates that the method can be very useful in cases where the animals' poses and viewpoints significantly vary and illumination conditions are highly different. However, this assessment is oracle-based since ground truth masks were employed.
\end{itemize}

In the multi-task framework, training of three tasks together proved to be highly conflicting.:
\begin{itemize}
    \item \textbf{Monolithic Hard-Shared Multi-Task Control (E1)}: The Monolithic Hard-Shared Multi-Task Control (E1) enforced all three tasks to share the same four-channel ResNet-18 spatial backbone. The performance of all three tasks was degraded compared to the single-task models. For BCS, MAE rose from 0.1709 to 0.1788. For Behavior, balanced accuracy was reduced from 74.43\% to 67.30\% and Walking F1 was down from 0.2456 to 0.0408. For Re-ID, Barn Rank-1 was reduced from 63.90\% to 57.38\% and mAP fell from 40.68\% to 30.37\%.
    \item \textbf{Modular Task-Private Multi-Task Model (E3)}: Modular Task-Private Multi-Task Learning (E3) added residual bottleneck adapters for each task between the shared backbone and task heads. It introduced 395,904 additional trainable parameters into the model, which is 3.32\% higher compared to the capacity of E1. However, it was unable to avoid negative transfer on a general level. The loss of behavior test decreased to 0.4193 and accuracy of CVB surveillance increased to 80.33\% in comparison to E1. At the same time, BCS became worse with MAE of 0.1916, and the Re-ID task performance became lower -- 49.08\% of Barn Rank-1 and 28.04\% of mAP. The walking F1 score was equal to 0.0000 on the matched held-out test.
    \item \textbf{Gradient-Projected Hard-Shared Multi-Task Optimization Control (E4)}: Gradient-Projected Hard-Shared Multi-Task Optimization Control (E4) was similar to E1 in number of 11,926,706 parameters. In addition, E4 applied PCGrad technique~\cite{yu2020pcgrad} only to the shared backbone during training. In case of E4 optimization, the number of projections with gradient conflicts is equal to 44,177 in 16,140 super-steps, wherein each super-step has 2.737 projections on average. It means that gradient conflicts happen quite often when training task pairs with the shared backbone in which conflict rates range from 46.5\% to 47.8\%. As for PCGrad comparison with E1, the model demonstrates accuracy equal to 86.92\%, balanced accuracy equal to 69.15\%, Macro-F1 equal to 0.7114, and walking F1 equal to 0.0909. Moreover, it managed to recover the important BCS and Re-ID results in contrast with E3. Still, E4 underperforms in comparison with the single-task models concerning some important metrics.
\end{itemize}

\section{Answers to the Research Questions}

The empirical findings from the single-task and multi-task investigations provide direct answers to the three research questions established in Chapter~1.

\subsection{Research Question 1}
\textit{How do cattle-centered representations affect BCS, behavior recognition, and Re-ID relative to their RGB single-task references?}

As for the effect of cattle-centered representations compared to the RGB single-task models, the enhancement of the key prediction and retrieval results was achieved in all three tasks. Yet, the effect varies depending on the task and metric combination.:
\begin{itemize}
    \item As for BCS, cattle-centered cropping and foreground mask help reduce the value of MAE by 0.0220 BCS units and improve Acc@1 by 4.45 percentage points in the matched test samples. The most significant improvements occurred in ordered error and tolerance metrics, while no improvements could be observed in the class-balanced metrics.
    \item As for Behavior Recognition, introducing the information about the foreground with temporal processing affected the results in terms of various classes and camera sources. The balanced accuracy has been improved on 3.13 percentage points, test loss decreased by 19.5\%, and Walking F1 score has become 0.0282 higher. On the other hand, the overall accuracy became slightly lower by 1.02 percentage points, and the accuracy for CVB surveillance decreased as well. Thus, the use of the temporal foreground approach could be beneficial for some movement-related tasks, though there was no equal improvement in all the metrics.
    \item For Re-ID, the target-centered representation worked best under conditions when images had high variations. The Snapshot-to-Parlor Rank-1 has improved on 24.05 percentage points, and mAP metric on 13.37 percentage points. As for the barn surveillance case, Rank-1 has become better only on 5.26 percentage points.
\end{itemize}
Thus, the cattle-centered representations were useful, yet not equally effective for all the metrics, and they did not completely solve the problem of shortcut learning between the tasks.

The improvements mentioned above were made possible thanks to the application of more than one processing layer. Both Cropping and Masking layers were used in BCS and Re-ID, while Behavior utilized both of those in addition to Temporal Convolution layer. Thus, the improvement could not have been made thanks to segmentation alone.

\subsection{Research Question 2}
\textit{How can temporal information and task-specific visual cues be incorporated while preserving the different requirements of condition assessment, activity recognition, and identity matching?}

This paper also analyzed how temporal features and visual cues can be integrated in consideration of diverse needs of condition scoring, behavior recognition and identity verification tasks.:
\begin{itemize}
    \item For Body Condition Scoring, body shape is the key biological information. The use of bounding-box crops and binary foreground masks made it easier for the model to focus on the cow while retaining the biological information of body shape that was essential for BCS. This improved prediction error without using information from multiple frames.
    \item In the case of Behavior Recognition, body pose and movement are both important aspects, thus requiring data from multiple frames for the model. In fact, a model using 8 frames with foreground guided TCN performed better than the single-frame model.
    \item For Cattle Re-ID system, details of the coat patterns and skin texture are required. The use of crops and segmentation masks played an important role in this regard. The Oracle masks turned out to be extremely useful in difficult conditions where the correct segmentation of the foreground was done.
    \item Anatomical pose keypoints and viewpoints have been thoroughly tested. The frozen SuperAnimal pose had a very poor detection rate (45.60\% in behavior sequence cases and 72.92\% in lying cattle cases) and was unable to generate the rear pelvic keypoints required for BCS. For these reasons, the pose information is not incorporated into the primary downstream pipeline. The real-cattle viewpoint had an 86.26\% accuracy rate on the dataset created just for it, but it has not been incorporated into the primary models since it has not been tested in other farm environments.
\end{itemize}
Thus, visual information has to be chosen based on requirements for each specific task. The body shape and morphology are needed for BCS, for behavior recognition change over time is necessary, and unique coat surface details are necessary for re-identification. Thus, it would be more useful to do that instead of using the same visual information for all the tasks.

\subsection{Research Question 3}
\textit{How does task-conditioned sharing compare with hard sharing in task performance, negative transfer, and robustness under domain change?}

Moreover, it was tested how well task-conditioned sharing and hard sharing affect the performance of tasks, negative transfer and robustness to domain shift. Comparing multi-task models with corresponding single-task models helped us make four important conclusions.:
\begin{itemize}
    \item \textbf{Outcome-Level Negative Transfer under Hard Sharing}: Hard parameter sharing model (E1) performed worse than corresponding single-task models in all three tasks. In BCS, MAE rose by 0.0079. For Behavior, the balanced accuracy dropped by 7.13 percentage points while the F1 score for Walking fell to 0.0408. In Re-ID, Barn Rank-1 declined by 6.52 percentage points, while mAP decreased by 10.31 percentage points. From these findings, we note that there is task outcome negative transfer whereby different cattle vision tasks were restricted to use the same ResNet-18 spatial backbone.
    \item \textbf{Evaluation of Task-Conditioned Modularity}: The introduction of task-specific residual bottleneck adapters in E3 had little impact on addressing the negative transfer problem. When compared to E1, E3 helped in reducing test loss for Behavior from 0.6226 to 0.4193 and CVB surveillance accuracy increased from 76.78\% to 80.33\%. Nevertheless, BCS MAE rose to 0.1916, Barn Rank-1 in Re-ID dropped to 49.08\%, and Walking F1 decreased to 0.0000. In addition, E3 gained 395,904 trainable parameters, hence increasing its model capacity by 3.32\% compared to E1. Due to this, the difference in performance cannot be justified purely based on the task-specific routing.
    \item \textbf{Evaluation of Gradient-Projected Optimization}: The evaluation of gradient-projected optimization was done using E4 where it applied PCGrad and had exactly the same number of parameters as E1 without adding additional capacity for the model. Conflicts between the gradients of task pairs occurred regularly throughout the process of training. In total, there were 44,177 conflict projections, out of 16,140 super-steps, where each super-step had an average of 2.737 projections and conflict rates ranging from 46.5\% to 47.8\%. PCGrad projected conflicting gradient components during optimization and provided selective performance mitigation. Behavior accuracy improved to 86.92\%, an increase of 1.92 percentage points from E1. Balanced accuracy improved to 69.15\%, an increase of 1.85 percentage points from E1, Macro-F1 improved to 0.7114 with an improvement of 0.0248, and Walking F1 improved to 0.0909 from 0.0408 in E1. It also recovered important BCS and Re-ID results compared to E3. In addition, E4 was unable to achieve full negative transfer avoidance. Single task models were superior in many ways. For BCS, MAE was 0.1709 compared to 0.1828 in E4. For Behavior, balanced accuracy was 74.43\% compared to 69.15\%, while walking F1 was 0.2456 compared to 0.0909. Finally, for Re-ID, Barn Rank-1 was 63.90\% compared to 54.97\%.
    \item \textbf{Robustness under Domain Change}: The multi-task models were evaluated under various circumstances, such as Barn and Snapshot views in SideView and CVB surveillance compared to Kaggle Beef feedlot videos for Behavior. Nevertheless, these evaluations only occurred using the datasets utilized in this thesis and thus do not confirm any robustness to other external commercial herds.
\end{itemize}

\section{Contributions of the Thesis}

This thesis makes six contributions to cattle monitoring and multi-task representation learning.:
\begin{enumerate}
    \item \textbf{Evaluation Protocols}: Several exact, reliable, and appropriate protocols have been developed to evaluate cattle monitoring tasks during the study process. The burst-group-disjoint protocol was adopted in the ScienceDB BCS case to prevent data leakage from the same passage frames. In Behavior Recognition using CVB and Kaggle Beef, data splitting was performed at both session and source levels. An identity-disjoint open-set protocol was employed in SideView Protocol A of Re-ID to validate 69 novel cows against a gallery of 36,811 parlor images.
    \item \textbf{Single-Task Baselines on Standard Dataset}: Normal RGB inputs versus cattle-centered representation were compared in the study using the same test samples. This comparison highlighted the benefits that come with the use of cattle cropping and foreground segmentation, providing substantial but task-dependent performance enhancement. For instance, MAE of BCS reduced from 0.1929 to 0.1709 and Snapshot Rank-1 of Re-ID increased from 38.88\% to 62.93\%.
    \item \textbf{Task-Specific Cattle Input Representation}: The study suggested a task-specific approach whereby input representations should be designed based on task requirements. For BCS, static cow crops and masks were used for capturing body shape information. For Behavior, 1D temporal convolution was applied to an 8-frame cow video sequence. For Re-ID, essential surface patterns were retained in the images. The limitations in terms of pose estimation and viewpoint were also discussed.
    \item \textbf{Benchmarking for Systematic Multi-Task Sharing}: In the systematic multi-task benchmark experiment, hard parameter sharing in E1, task-private residual adapter in E3, and gradient-projected optimization in E4 were compared to their respective single-task baselines. The same held-out test sets were used in the benchmark to evaluate the multi-task cattle vision system.
    \item \textbf{Direct Gradient Conflict Diagnostics}: In this study, direct evidence of gradient conflicts in multi-task learning with PCGrad was presented. A total of 44,177 gradient conflicts were found in 16,140 super-steps. The rate of conflict is between 46.5\% and 47.8\%. This implies that conflicts occur regularly during multi-task cattle monitoring training.
    \item \textbf{Multi-Task Trade-Offs}: As per the research findings, multi-task learning is not necessarily advantageous to all cattle monitoring tasks. In case tasks have shared features, one task will be improved while the other gets impaired. Negative transfer refers to the scenario where one solution improves one metric but not others.
\end{enumerate}

\section{Limitations}

A number of limitations in the study need to be explicitly stated.:
\begin{enumerate}
    \item \textbf{Absence of Verified Cow IDs in ScienceDB}: ScienceDB does not verify the IDs of cows. Though, the burst group disjoint split does not allow very similar frames from the same passage to appear in different sets, it does not guarantee that different sets have different cows.
    \item \textbf{Session-Disjoint Rather than Cow-Disjoint Behavior Evaluation}: For Behavior task, the CVB and Kaggle Beef data have been split by videos and sessions rather than cows. Hence, Behavior task evaluation is session disjoint and not cow disjoint.
    \item \textbf{Confounding of Walking class}: The Walking class was observed solely in the CVB surveillance dataset, where there were 26 test sequences. Therefore, there is some association between the camera source and the Walking action that could affect the classification performance.
    \item \textbf{Automatic Perception Pipeline Coverage}: The coverage of the BCS perception pipeline was 93.89\% with respect to the ScienceDB test set. In 8,040 test images, there were 489 detection failures and 2 segmentation failures. The Behavior pipeline retained 96.42\% of the test sequences. Images with perception failure were excluded from our evaluation.
    \item \textbf{Guidance from Oracle Segmentation in Re-ID}: In Run 6 and other multi-task Re-ID experiments, we relied on the oracle masks available in the dataset rather than on the output of a segmentation model. This provides an upper bound of the foreground isolation result, however, it does not reflect an automatic end-to-end solution.
    \item \textbf{Joint Component Attribution}: The single-task models in Runs 4--6 were composed using different components like localization, cropping, binary mask, and temporal convolutions altogether. Hence, the contribution of each component was not measured individually.
    \item \textbf{Confounding in Capacity of Modular Architecture}: E3 introduced 395,904 additional trainable parameters, increasing the model size by 3.32\% relative to E1. For this reason, the performance difference between E1 and E3 is not only due to routing but because the architecture was also bigger.
    \item \textbf{Boundary on Causal Explanation for Gradient Conflict}: Gradient conflict has been directly observed only when training E4 using PCGrad. The experiment shows the presence of gradient conflicts during E4 training. However, this is not sufficient to assert gradient conflict being the sole cause of worse performance in E1 as well as E1 having similar gradient conflict rate as E4.
    \item \textbf{Point-Estimate Evaluation}: All single-task and multi-task models have been evaluated using only one checkpoint from training. Since computational limitations restricted our ability to run the training multiple times with different seeds, confidence intervals and hypotheses testing could not be conducted. Hence, all of the observed differences between the metrics are merely descriptive point comparisons.
    \item \textbf{Untested Experiments in External Stress Testing Domain}: While we conducted an evaluation on the benchmark datasets, an external stress test on other types of datasets such as Ruchay, Dryad, MmCows, CBVD-5, and BECA datasets had been planned but was not executed. Hence, the obtained results cannot be considered for any verification of the cross-domain generalization.
    \item \textbf{Unverified Downstream Transfer of Pose and Viewpoint}: While conducting the feasibility tests, both pose estimation and viewpoint were highly prone to failures as well as learning some source-specific shortcuts. Hence, the usefulness of them for downstream task datasets was not validated.
    \item \textbf{Negative Transfer Persisting after Improvement of Results with PCGrad}: Although some improvement was achieved with PCGrad, the performance of multi-task models remained inferior compared to the single-task models on many key metrics, showing that joint multi-task deployment requires an engineering trade-off.
\end{enumerate}

\section{Future Work}

There are several potential research directions that could be pursued in the future in light of this paper's results and limitations.:
\begin{itemize}
    \item \textbf{Multiple Seed Training and Uncertainty Estimates}: In future research, it would be useful to repeat the training procedure with different seeds in order to obtain estimates for confidence intervals, variance and statistical testing of both one and multi-task models.
    \item \textbf{End-to-end Auto-Re-ID Segmentation}: In further research, an automatic segmentation model, e.g., fine-tuned SAM or YOLO-seg, could be used to generate prediction masks instead of oracle masks. It might help us to investigate how the errors made by the segmentation network affect the Re-ID task in realistic settings.
    \item \textbf{Generalization via External Testing with Multiple Herds}: In addition, the trained models should also be tested against independent external datasets. In case of BCS, it could be Ruchay2026 and Dryad, whereas in case of Behavior Recognition, MmCows and CBVD-5 would do.
    \item \textbf{Cow Identity Retrieval over Time}: One possible approach for doing future work may be through cow Re-ID over time via the datasets BECA-L and BECA-D. It would be able to reveal the stability of cow identity attributes across different groupings, lactations, season coat, and body development phases.
    \item \textbf{Ablation Study of Individual Perception Components}: An ablation study for the effect of bounding box, soft mask, binary foreground mask, and sequence lengths could be performed.
    \item \textbf{Robust Pose Estimation of Animal}: Fine-tuning of animal pose estimation models for the laying and rear view of the cattle could be carried out.
    \item \textbf{Multi-task Optimization Approaches}: One can perform future experiments using partial parameter sharing approaches like branched trunk and adaptive loss weighting approaches like GradNorm~\cite{chen2018gradnorm} to resolve the conflict of gradients between different tasks.
\end{itemize}

\clearpage
\section{Conclusion}

The thesis was successful in establishing an empirical and methodological framework for multi-task deep learning in the context of precision cattle monitoring, involving body condition scoring, behavior recognition, and cow re-identification under controlled single-task and joint multi-task settings.

Evidence from experiments has shown that cow-focused visual representations offer significant utility in the performance of specific individual livestock monitoring tasks by improving BCS error tolerance and improving the re-identification of cows when there are unconstrained shifts in the cameras. This is however task dependent and does not address the separate needs of representation for morphological regression, temporal modeling, and biometric identification.

Naive hard parameter sharing in a multi-task integration framework led to negative transfer effects at the outcome level in all three tasks. Adding task-private residual adapters did not help mitigate the performance drop, whereas PCGrad explicitly showed that gradient directions for shared backbones were often opposite in training. Through the balancing of opposing gradient components, PCGrad delivered a selective remedy: especially improving behavior classification metrics and CVB surveillance performance, and yet specific single-task models were better on some task-specific metrics.

In the end, the results show that multi-task learning in agricultural computer vision is not universally beneficial in terms of performance using shared representation; successful deployment should entail a balance between shared characteristics of cattle, task-oriented abilities, and optimization strategies that account for task conflicts.

